% Options for packages loaded elsewhere
% Options for packages loaded elsewhere
\PassOptionsToPackage{unicode}{hyperref}
\PassOptionsToPackage{hyphens}{url}
%
\documentclass[
  ignorenonframetext,
  aspectratio=169,
  russian,
]{beamer}
\newif\ifbibliography
\usepackage{pgfpages}
\setbeamertemplate{caption}[numbered]
\setbeamertemplate{caption label separator}{: }
\setbeamercolor{caption name}{fg=normal text.fg}
\beamertemplatenavigationsymbolshorizontal
% Prevent slide breaks in the middle of a paragraph
\widowpenalties 1 10000
\raggedbottom
\AtBeginPart{
  \frame{\partpage}
}
\AtBeginSection{
  \ifbibliography
  \else
    \frame{\sectionpage}
  \fi
}
\AtBeginSubsection{
  \frame{\subsectionpage}
}
\usepackage{iftex}
\ifPDFTeX
  \usepackage[T1]{fontenc}
  \usepackage[utf8]{inputenc}
  \usepackage{textcomp} % provide euro and other symbols
\else % if luatex or xetex
  \usepackage{unicode-math} % this also loads fontspec
  \defaultfontfeatures{Scale=MatchLowercase}
  \defaultfontfeatures[\rmfamily]{Ligatures=TeX,Scale=1}
\fi
\usepackage{lmodern}

\usefonttheme{serif} % use mainfont rather than sansfont for slide text
\ifPDFTeX\else
  % xetex/luatex font selection
  \setmainfont[]{Liberation Serif}
  \setsansfont[]{Liberation Sans}
  \setmonofont[]{Liberation Mono}
\fi
% Use upquote if available, for straight quotes in verbatim environments
\IfFileExists{upquote.sty}{\usepackage{upquote}}{}
\IfFileExists{microtype.sty}{% use microtype if available
  \usepackage[]{microtype}
  \UseMicrotypeSet[protrusion]{basicmath} % disable protrusion for tt fonts
}{}


\usepackage{longtable,booktabs,array}
\usepackage{calc} % for calculating minipage widths
\usepackage{caption}
% Make caption package work with longtable
\makeatletter
\def\fnum@table{\tablename~\thetable}
\makeatother
\usepackage{graphicx}
\makeatletter
\newsavebox\pandoc@box
\newcommand*\pandocbounded[1]{% scales image to fit in text height/width
  \sbox\pandoc@box{#1}%
  \Gscale@div\@tempa{\textheight}{\dimexpr\ht\pandoc@box+\dp\pandoc@box\relax}%
  \Gscale@div\@tempb{\linewidth}{\wd\pandoc@box}%
  \ifdim\@tempb\p@<\@tempa\p@\let\@tempa\@tempb\fi% select the smaller of both
  \ifdim\@tempa\p@<\p@\scalebox{\@tempa}{\usebox\pandoc@box}%
  \else\usebox{\pandoc@box}%
  \fi%
}
% Set default figure placement to htbp
\def\fps@figure{htbp}
\makeatother



\ifLuaTeX
\usepackage[bidi=basic,provide=*]{babel}
\else
\usepackage[bidi=default,provide=*]{babel}
\fi
\ifPDFTeX
\else
\babelfont{rm}[]{Liberation Serif}
\fi
% get rid of language-specific shorthands (see #6817):
\let\LanguageShortHands\languageshorthands
\def\languageshorthands#1{}


\setlength{\emergencystretch}{3em} % prevent overfull lines

\providecommand{\tightlist}{%
  \setlength{\itemsep}{0pt}\setlength{\parskip}{0pt}}



 

\usepackage[]{csquotes}

\IfFileExists{plex-otf.sty}{
  %% Full TeXlive
  % \usepackage[%
  %   % math,
  %   RM={Scale=0.94},SS={Scale=0.94},SScon={Scale=0.94},TT={Scale=MatchLowercase,FakeStretch=0.9},DefaultFeatures={Ligatures=Common}
  % ]{plex-otf}
}{
  %% TinyTeX
  \usepackage{libertine}
}

%%% Load theme
% https://deic.uab.cat/~iblanes/beamer_gallery/
\IfFileExists{beamerthemegotham.sty}{
  %% Full TeXlive
  \usetheme{gotham}
  \gothamset{
    numbering=totalpagenumber,
    parttocframe default=off,
    sectiontocframe default=off,
    subsectiontocframe default=off,
  }
}{
  %% TinyTeX
  \usetheme{Madrid}
}
\makeatletter
\@ifpackageloaded{caption}{}{\usepackage{caption}}
\AtBeginDocument{%
\ifdefined\contentsname
  \renewcommand*\contentsname{Содержание}
\else
  \newcommand\contentsname{Содержание}
\fi
\ifdefined\listfigurename
  \renewcommand*\listfigurename{Список иллюстраций}
\else
  \newcommand\listfigurename{Список иллюстраций}
\fi
\ifdefined\listtablename
  \renewcommand*\listtablename{Список таблиц}
\else
  \newcommand\listtablename{Список таблиц}
\fi
\ifdefined\figurename
  \renewcommand*\figurename{Рисунок}
\else
  \newcommand\figurename{Рисунок}
\fi
\ifdefined\tablename
  \renewcommand*\tablename{Таблица}
\else
  \newcommand\tablename{Таблица}
\fi
}
\@ifpackageloaded{float}{}{\usepackage{float}}
\floatstyle{ruled}
\@ifundefined{c@chapter}{\newfloat{codelisting}{h}{lop}}{\newfloat{codelisting}{h}{lop}[chapter]}
\floatname{codelisting}{Список}
\newcommand*\listoflistings{\listof{codelisting}{Листинги}}
\makeatother
\makeatletter
\makeatother
\makeatletter
\@ifpackageloaded{caption}{}{\usepackage{caption}}
\@ifpackageloaded{subcaption}{}{\usepackage{subcaption}}
\makeatother

\usepackage{bookmark}
\IfFileExists{xurl.sty}{\usepackage{xurl}}{} % add URL line breaks if available
\urlstyle{same}
\hypersetup{
  pdftitle={Работа с репозиториями Git и GitFlow},
  pdfauthor={Никифоров Захар Сергеевич},
  pdflang={ru-RU},
  hidelinks,
  pdfcreator={LaTeX via pandoc}}


\title{Работа с репозиториями Git и GitFlow}
\subtitle{Лабораторная работа №4}
\author{Никифоров Захар Сергеевич}
\date{2026-03-19}

\begin{document}
\frame{\titlepage}

\renewcommand*\contentsname{Содержание}
\begin{frame}[allowframebreaks]
  \frametitle{Содержание}
  \setcounter{tocdepth}{2}
  \tableofcontents
\end{frame}
\setcounter{tocdepth}{2}
\tableofcontents
}

\section{1. Вводная
часть}\label{ux432ux432ux43eux434ux43dux430ux44f-ux447ux430ux441ux442ux44c}

\begin{frame}{1.1 Актуальность}
\phantomsection\label{ux430ux43aux442ux443ux430ux43bux44cux43dux43eux441ux442ux44c}
\begin{itemize}
\tightlist
\item
  Git --- основной инструмент для управления версиями кода
\item
  GitFlow --- позволяет соблюдать структуру в разработке
\item
  Знание правильной работы с репозиториями
\end{itemize}
\end{frame}

\begin{frame}{1.2 Объект и предмет исследования}
\phantomsection\label{ux43eux431ux44aux435ux43aux442-ux438-ux43fux440ux435ux434ux43cux435ux442-ux438ux441ux441ux43bux435ux434ux43eux432ux430ux43dux438ux44f}
\begin{itemize}
\tightlist
\item
  Объект исследования: процесс работы с репозиториями Git
\item
  Предмет исследования: GitFlow и инструменты работы с репозиториями
\end{itemize}
\end{frame}

\begin{frame}{1.3 Цели и задачи}
\phantomsection\label{ux446ux435ux43bux438-ux438-ux437ux430ux434ux430ux447ux438}
Цель: Получение навыков правильной работы с репозиториями git.

Задачи: 1. Установить необходимое ПО 2. Создать репозиторий git-extended
3. Выполнить его синхронизацию с GitHub 4. Настроить конфигурацию
коммитов 5. Ознакомится с системой релизов
\end{frame}

\begin{frame}{1.4 Материалы и методы}
\phantomsection\label{ux43cux430ux442ux435ux440ux438ux430ux43bux44b-ux438-ux43cux435ux442ux43eux434ux44b}
\begin{itemize}
\tightlist
\item
  Оборудование: ПК с ОС Linux или ВМ на нем с ОС Linux
\item
  ПО: git, gh, git-flow, node.js, pnpm
\item
  Методы: Установка согласно руководству, работа с ветками, создание
  релизов.
\end{itemize}
\end{frame}

\section{2. Порядок выполнения
работы}\label{ux43fux43eux440ux44fux434ux43eux43a-ux432ux44bux43fux43eux43bux43dux435ux43dux438ux44f-ux440ux430ux431ux43eux442ux44b}

\begin{frame}{2.1 Установка программ}
\phantomsection\label{ux443ux441ux442ux430ux43dux43eux432ux43aux430-ux43fux440ux43eux433ux440ux430ux43cux43c}
\begin{itemize}
\tightlist
\item
  Устанавливаем git, gh, git-flow, node.js, pnpm
\end{itemize}

\begin{figure}

\centering{

\includegraphics[width=0.7\linewidth,height=\textheight,keepaspectratio]{image/1.png}

}

\caption{\label{fig-001}Скриншот терминала 1}

\end{figure}%

\begin{figure}

\centering{

\includegraphics[width=0.7\linewidth,height=\textheight,keepaspectratio]{image/2.png}

}

\caption{\label{fig-002}Скриншот терминала 2}

\end{figure}%
\end{frame}

\begin{frame}{2.2 Создание Git-репозитория}
\phantomsection\label{ux441ux43eux437ux434ux430ux43dux438ux435-git-ux440ux435ux43fux43eux437ux438ux442ux43eux440ux438ux44f}
\begin{itemize}
\tightlist
\item
  Создаем локальный репозиторий и инициализируем, получаем первый
  README.md, который уже закомментировали.
\end{itemize}

\begin{figure}

\centering{

\includegraphics[width=0.7\linewidth,height=\textheight,keepaspectratio]{image/3.png}

}

\caption{\label{fig-003}Скриншот терминала 3}

\end{figure}%
\end{frame}

\begin{frame}{2.3 Настройка commitizen}
\phantomsection\label{ux43dux430ux441ux442ux440ux43eux439ux43aux430-commitizen}
\begin{itemize}
\tightlist
\item
  Используя утилиту commitizen, делаем коммит с помощью git cz
\end{itemize}

\begin{figure}

\centering{

\includegraphics[width=0.7\linewidth,height=\textheight,keepaspectratio]{image/4.png}

}

\caption{\label{fig-004}Скриншот терминала 4}

\end{figure}%
\end{frame}

\begin{frame}{2.4 Настройка файла package.json}
\phantomsection\label{ux43dux430ux441ux442ux440ux43eux439ux43aux430-ux444ux430ux439ux43bux430-package.json}
\begin{itemize}
\tightlist
\item
  Открываем package.json и редактируем его.
\end{itemize}

\begin{figure}

\centering{

\includegraphics[width=0.7\linewidth,height=\textheight,keepaspectratio]{image/05.png}

}

\caption{\label{fig-005}\enquote*{package.json}}

\end{figure}%
\end{frame}

\begin{frame}{2.5 Инициализация Git Flow и создание ветки feature.}
\phantomsection\label{ux438ux43dux438ux446ux438ux430ux43bux438ux437ux430ux446ux438ux44f-git-flow-ux438-ux441ux43eux437ux434ux430ux43dux438ux435-ux432ux435ux442ux43aux438-feature.}
\begin{itemize}
\tightlist
\item
  Инициализируем Git Flow
\item
  Получаем базовые ветки master и develop
\item
  Создаем ветку feature
\end{itemize}

\begin{figure}

\centering{

\includegraphics[width=0.7\linewidth,height=\textheight,keepaspectratio]{image/8.png}

}

\caption{\label{fig-006}Скриншот терминала 5}

\end{figure}%
\end{frame}

\begin{frame}{2.6 Проверка веток и создание журнала изменений}
\phantomsection\label{ux43fux440ux43eux432ux435ux440ux43aux430-ux432ux435ux442ux43eux43a-ux438-ux441ux43eux437ux434ux430ux43dux438ux435-ux436ux443ux440ux43dux430ux43bux430-ux438ux437ux43cux435ux43dux435ux43dux438ux439}
\begin{itemize}
\tightlist
\item
  По завершении работы с веткой feature, финишируем ее
\item
  Проверяем ветки
\item
  Создаем новый релиз 1.2.3 и завершаем его.
\end{itemize}

\begin{figure}

\centering{

\includegraphics[width=0.7\linewidth,height=\textheight,keepaspectratio]{image/9.png}

}

\caption{\label{fig-007}Скриншот терминала 6}

\end{figure}%
\end{frame}

\section{3. Выводы}\label{ux432ux44bux432ux43eux434ux44b}

\begin{frame}{3. Выводы}
\begin{itemize}
\tightlist
\item
  Были получены навыки правильной работы с репозиториями git.
\end{itemize}
\end{frame}




\end{document}
